% Research note for the parent agent. No external logic or coverage is used.
\section{Predecessor without numerical coverage}
\label{sec:equational-predecessor}

Let $\mathcal E$ be the purely equational calculus on $0,S,I$, with
ordinary equality rules, the two iterator equations, and literal-equation
induction with parameters.  In particular, $\mathcal E$ has no axiom saying
that every number is zero or a successor.  The following strengthens the
dependency statement of the inherited scalar predecessor construction:
\[
  \mathcal E\vdash P(0)=0,
  \qquad \mathcal E\vdash P(Sx)=x
\]
for one finite iterator term $P$.  Thus $\mathcal E$ proves successor
injectivity.  The construction below is inherited; the new point is the
explicit elimination of all its numerical and Boolean case splits.

Write $1=S0$ and define
\[
 C(e,a,b)=I(e,a;u.b),\qquad Z(e)=C(e,1,0),\qquad
 a+b=I(b,a;u.Su).
\]
The equations of a commutative monoid for addition, and
$Sa+b=S(a+b)$, follow by equation induction.  Expressions such as $2a$ and
$3a$ abbreviate repeated addition.

\paragraph{Equational interpolation.}
For every term template $F(v)$, with all other variables fixed, equation
induction on a fresh $e$ proves
\begin{equation}
 F(Z(e))=C(e,F(1),F(0)). \tag{BI}
\end{equation}
At $e=0$ both sides are $F(1)$; at $e=Sz$ both sides are $F(0)$.
The induction hypothesis is unused.  Also equation induction, again with
an unused hypothesis, proves $Z(Z(Z(e)))=Z(e)$.
Consequently, if $b=Z(Z(b))$, and the two equations
$F(0)=G(0)$ and $F(1)=G(1)$ are available under the current fixed
parameters, then $F(b)=G(b)$ follows by applying (BI) at $e=Z(b)$.
This is a derived equality rule, not a disjunction elimination rule.  It
continues to hold in the presence of arbitrary fixed side equations.

\paragraph{Parity and a three-cycle.}
Put
\[
 N(n)=I(n,0;u.Z(u)),\quad j(u)=Z(u)+2N(u),\quad
 M(n)=I(n,0;u.j(u)).
\]
Equation induction proves $N(n)=Z(Z(N(n)))$: the zero case computes;
the successor case is $Z(N(n))=Z(Z(Z(N(n))))$.  Applying (BI) therefore
gives
\[
 N(n)+Z(N(n))=1,\qquad N(n+2)=N(n).
\]
Equation inductions give
\[
 N(2n)=0,\qquad N(2n+1)=1,\qquad N(3n)=N(n).
\]
For the last equation use $N(a+3)=Z(N(a))$ in the successor step.
Finite computations yield $j(0)=1,j(1)=2,j(2)=0$.

\paragraph{The ramp invariant, with its case split removed.}
Define
\[
 G(n)=I(n,0;u.u+M(u)+1).
\]
We prove $M(G(n))=N(n)$ by equation induction.  At the successor step
write $a=G(n)$ and $b=N(n)$, so the sole induction equation is $M(a)=b$.
Interpolation in the Boolean parameter $b$ gives
\[
 M(a+b+1)=C(b,M(a+1),M(a+2)).
\]
Both endpoint equations follow from the computations of $C$.  Rewriting
by $M(a)=b$ and the two computations of $M$ turns the right side into
$C(b,j(b),j(j(b)))$.  A further use of (BI) gives
\[
 C(b,j(b),j(j(b)))=Z(b):
\]
at $b=0$ both sides are $1$, and at $b=1$ both sides are $0$.
Since $G(Sn)=a+b+1$ and $N(Sn)=Z(b)$, this proves the induction step.
No disjunctive fact about $b$ was used.

It follows by equality reasoning and equation induction that
\[
 G(Sn)=G(n)+N(n)+1,\quad G(n+2)=G(n)+3,\quad G(2n)=3n.
\]
Another equation induction proves
\[
 G(3n)=3G(n)+N(n).
\]
Its successor step reduces, with $b=N(n)$, to the Boolean identity
$2b+4=3b+3+Z(b)$, whose endpoints are checked directly and whose
general Boolean instance follows by (BI).

\paragraph{The factor map.}
Set $H(u)=C(N(u),G(u),u)$.  We obtain
\[
 H(2u)=3u,\qquad H(1)=1,\qquad H(3u)=3H(u).
\]
Only the last identity requires comment.  With $b=N(u)$ it is
\[
 C(b,3G(u)+b,3u)=3C(b,G(u),u),
\]
which is again (BI): its zero endpoint is $3G(u)=3G(u)$ and its one
endpoint is $3u=3u$.  The apparent dependence between $b$ and $u$
causes no problem: prove the displayed identity with an independent
Boolean variable, then substitute $N(u)$.

\paragraph{Comparator.}
For any fixed update $h$, equation induction gives rotation
$I(Sy,a;h)=I(y,h(a);h)$.  Equation induction also gives
\[
 I(y,3a;H)=3I(y,a;H),\qquad I(y,1;H)=1.
\]
Let
\[
 E(x)=I(x,1;u.2u),\qquad L(x,y)=N(I(y,E(x);H)).
\]
The preceding identities give the complete comparator specification
\[
 L(0,y)=1,\qquad L(Sx,0)=0,\qquad L(Sx,Sy)=L(x,y).
\]
For the third equation the calculation is
\[
\begin{aligned}
 N(I(Sy,2E(x);H))
 &=N(I(y,H(2E(x));H))\\
 &=N(I(y,3E(x);H))\\
 &=N(3I(y,E(x);H))=N(I(y,E(x);H)).
\end{aligned}
\]

\paragraph{Saturation and predecessor.}
Put
\[
 B_m(u)=C(L(Su,m),u,Su),\qquad W(t,m)=I(t,0;u.B_m(u)).
\]
The identity $C(e,Sa,Sb)=SC(e,a,b)$ is an equation induction on $e$
with unused induction hypothesis.  It gives
\[
 B_{Sm}(Su)=SB_m(u),\qquad B_{Sm}(0)=S0.
\]
Equation induction on $t$, with $m$ fixed, gives
$W(St,Sm)=SW(t,m)$.  Equation inductions then give
\[
 L(Sn,n)=0,\qquad W(n,n)=n,\qquad W(Sn,n)=n.
\]
In the middle equation substitute $t=m=n$ into the already established
two-parameter shift equation.  No parameter of a live induction hypothesis
is changed.

Finally define the finite scalar term
\[
 P(m)=I(m,0;u.C(L(SSu,m),u,Su)).
\]
Its zero value computes.  The comparator recurrence identifies the update
of $P(Sm)$ with $B_m(u)$.  Formally, equation induction on a fresh counter
$t$ proves
\[
 I(t,0;u.C(L(SSu,Sm),u,Su))=W(t,m).
\]
At $t=Sm$ the right side is $m$, proving $P(Sm)=m$.  This last
induction accounts for iterator-body extensionality; it is not an extra
equality axiom.  Applying $P$ to $Sx=Sy$ proves $x=y$.

\paragraph{Audit boundary.}
The proof uses only the two iterator equations, equality reasoning,
substitution into established equations, and equation induction with fixed
side parameters.  Every Boolean step has been replaced by (BI).  It does
not use coverage, a predecessor axiom, successor injectivity, equality
reflection, pairing, or induction on a non-equational formula.  This is a
written syntactic derivation; it is not a new Lean verification.
